\section{示例库：框架究竟如何形成}

\subsection{示例 1：学习进入仪式}

\begin{example}[重复进入线索制造廉价的重新进入]
一个学生反复在同一地点开始工作，打开同一本笔记到下一个未解决的位置，并以同一个五分钟起步任务进入状态。
随着时间推移，在短暂中断之后再次进入学习状态会变得更容易。
\end{example}

\begin{discussion}
这是通过重复耦合实现涌现的一个干净案例。
局部事件很简单：同一地点、同一线索、同一进入动作。
高层结果则是一个学习不再需要从零重建的框架。
这里的证据车道是\textbf{中等}：复杂性与涌现教材，加上亚稳态综述，足以支撑解释结构，
但当前来源集中并没有强实证论文直接证明这一具体学习仪式。
\end{discussion}

\subsection{示例 2：刷手机作为一种干扰架构}

\begin{example}[竞争路径获得历史优先权]
设某人反复在每一个短暂空档期查看手机：写完一段之后、吃完饭之后、排队时、睡前。
随着时间推移，“空闲到手机”这一路径会比“空闲到恢复工作”更快、更可预期，也更能自我强化。
\end{example}

\begin{discussion}
这个例子展示的不是简单的习惯累积，而是干扰。
旧路径现在会侵入本来无关的情境，因为它获得了更优越的重新进入经济性。
用系统语言说，手机路径已经成为一个高度可达的相邻盆地。
其解释价值属于\textbf{中等}，来源于路径依赖、涌现与吸引子语言，而不是任何直接神经解码式主张。
\end{discussion}

\subsection{示例 3：通过多尺度组织形成运动身份}

\begin{example}[框架跨越身体、日程与社会时序而形成]
一位上班族开始在三个固定晚上去健身房，提前摆好衣物，沿同一条路线出行，并在同一时间与同一个搭档会合。
几周之后，“运动之夜”不再像每次都要重新做出的道德决定。
\end{example}

\begin{discussion}
这一案例很重要，因为它是多尺度的。
框架并不是仅靠内在决心形成的。
身体准备、路线重复、日历结构与社会同步都在共同作用。
这个宏观模式之所以具有解释价值，是因为它改变了这个夜晚\emph{是什么}。
系统现在携带的是一个机制，而不只是一个愿望。
\end{discussion}

\subsection{示例 4：研究写作与负向路径依赖}

\begin{example}[回避被组织成一种能力]
一名研究生反复通过“再读一篇论文、再记一条笔记、再润色一条引文”来处理不确定性，而不是开始写正文。
最终，写作前准备状态变得极其熟练，而真正的显性写作却越来越昂贵。
\end{example}

\begin{discussion}
这是负向涌现的一个好例子。
系统并不只是“没写出来”；它已经围绕避免暴露而完成了组织。
这一形成出来的框架具有解释价值，因为它改变了后续成本结构。
因此，手稿不应只问“为什么学生今天不愿意写？”，
还应问“哪一种反复执行的微观政策，把这种退缩训练成了一条稳定路径？”
\end{discussion}

\subsection{示例 5：群体规范形成}

\begin{example}[框架也可以是集体性的]
在一个团队会议中，人们反复奖励快速同意，并以隐性方式惩罚缓慢澄清。
足够多次重复之后，群体会形成一种框架：快速共识显得自然，而谨慎异议则显得具有社会成本。
\end{example}

\begin{discussion}
这个例子扩大了本书的适用范围。
框架形成并不只是个人层面的。
重复的局部互动可以生成集体运行模式，而这些模式会进一步支配个体感觉什么容易、什么昂贵。
复杂性教材与涌现理论在这里尤其有价值，因为这一宏观模式无法被还原为任何一个参与者的孤立意图。
\end{discussion}